\chapter*{Sammendrag}

Småsatellitteknologi har økt tilgjengeligheten til satellittdata for jordobservasjon. Satellittprosjektet HYPerspektral Småsatellitt for havObservasjon (\acrshort{hypso}) ved Norsk Teknisk-Naturvitenskapelig Universitet (NTNU) bruker hyperspektrale kamera med 120 frekvensbånd for observasjon av havet. For effektiv lagring og nedlasting av de store datamengdene produsert av kamera må bildekompresjon benyttes ombord i satellitten. Arbeidsgruppen '\acrlong{CCSDS}' (\acrshort{CCSDS}) utvikler algoritmer for kompresjon av hyperspektrale og multispektrale bilder. Av disse støtter algoritmen CCSDS 123.0-B-2 nær tapsfri kompresjon, som tillater høye kompresjonsrater. Kalkulasjonsavhengihet mellom sekvensielt prosesserte datapunkter skaper utfordringer for implementasjonen av kompresjons-akselleratorer for algoritmen. Dette arbeidet undersøker bruken av en ny metode for kalkulering av vekt-vektorer i algortimen, foreslått av \citeauthor{jiaRemovalFeedbackLoop2025}, for å øke prosesseringskapasiteten til akseleratoren utviklet av \citeauthor{vorhaugHyperspectralImageCompression2024}. For bilder tatt av \acrshort{hypso}-2 gir den nye metoden en reduksjon i kompresjonsrate på under $1\%$. Den modifiserte akseleratoren når en prosesseringskapasitet på 202 MS/s for \acrshort{FPGA}en Zynq ZU7EV, med det minste fotavtrykket blant kjente akseleratorer på 6349 LUTer og 3480 FFer. En kapasitet på 351 MS/s er nådd for akseleratorens prediksjon-modul, som gjør den til den raskeste i literaturen. Akseleratoren er testet ombord i satellitten \acrshort{hypso}-2 med en klokkefrekvens på 92 MHz på \acrshort{FPGA}en Zynq-7030.
